% !TeX root = ../nc-quadric-avatar.tex
% 2026-11-23 / scratch draft: vertex-ideal criterion and all classified avatars
\section{頂点idealによる反Serre生成判定と一般avatar}
\label{sec:qa-general-vertex-az}

前節の二捩れspinorによる障害を，一般のavatarへ無条件に
外挿することはできない．本節では先にAS helixの頂点idealに
よる生成判定を証明する．最小十本矢型では二本の追加矢が作る
無限pathを用いることで，二捩れを仮定せずに偶数頂点を排除できる．
その後，非被約曲線を許すsquare avatarについて，二つのmutation壁と
cubicのpointed逆構成を用いて全四位相を判定する．

\subsection{十分負の成分比較と頂点ideal}

$A$を本稿のsquare，rotated，cubic，または最小十本矢型の
左右Noetherian AS $\mathbb Z$-代数とする．
$P_i=e_iA$, $E_i=\pi P_i$，
$t=S^{-1}[2]$と書き，$tE_i\simeq E_{i+4}$の規約を用いる．
右加群の$j$成分は$M_j=Me_j$であり，
$(P_a)_j=A_{aj}$は$j\le a$に支持される．

\begin{lemma}[有限射影解消からの遠方成分比較]
\label{lem:qa-as-far-component-comparison}
任意の有限生成次数付き右$A$加群$M$に対し，$j\ll0$なら
自然な写像は同型
\[
 M_j\xrightarrow{\sim}\Hom_{\qgr A}(E_j,\pi M)
\]
を与え，さらに$\Ext^r_{\qgr A}(E_j,\pi M)=0$ ($r>0$)である．
必要な$j$の上限は$M$に依存してよい．
\end{lemma}
\begin{proof}
$A$は左右Noetherianであり，左右の単純加群は長さ三のAS解消を持つ．
従って$M$には有限個の標準射影$P_a$からなる長さ三以下の解消がある．
これは命題~\ref{prop:qa-as-avatar-coherent-morita}の証明と同じである：
右AS解消を双対化すると，AS--Gorenstein条件から全ての左単純の
長さ三の射影解消を得る．有限生成右加群は添字について上に有界で，
graded Nakayamaによる極小射影被覆を持つ．
有限生成syzygyの極小射影被覆を取ると，第四項が非零なら
いずれかの左単純との$\Tor_{>3}$で検出され，
左単純の長さ三の解消に矛盾する．
AS10についても定理~\ref{thm:qa-as10-global-bijection}および
命題~\ref{prop:qa-as10-noetherian-qgr}により同じ入力が成立する．

この有限解消を$P_\bullet\to M$とする．
補題~\ref{lem:qa-as-serre-unique-recovery}，
補題~\ref{lem:qa-as10-hilbert-foundation}および全整数mutationの
前向きExt消滅より，$j$を解消に現れる全$a$より十分小さく取れば
\[
 \RHom(E_j,\pi P_a)=A_{aj}[0]
\]
であり，同型は全写像と両立する．
従って$\RHom(E_j,\pi P_\bullet)$は元の射影解消の$j$成分そのもので，
そのコホモロジーは次数零の$M_j$だけである．
$\pi$は完全なので，これは$\RHom(E_j,\pi M)$を計算する．
ここでは任意の個別頂点の豊富性や，一般的な$\chi$定理を先に仮定していない．
\end{proof}

\begin{theorem}[頂点idealによる反Serre豊富性の必要十分条件]
\label{thm:qa-vertex-ideal-az-criterion}
$i\in\mathbb Z$を固定し，
\[
 I_i=\sum_{q\in\mathbb Z}Ae_{i+4q}A,
 \qquad C_i=A/I_i
\]
と置く．次は同値である：
\begin{enumerate}
\item $(E_i,t)$はArtin--Zhang ampleである．
\item すべての$j\in\mathbb Z$について右$C_i$加群$e_jC_i$は
      有限次元である．
\end{enumerate}
特に，削除後に無限次元の標準右射影が残る位相は
強い$\HH$の構造頂点にならない．
\end{theorem}
\begin{proof}
(2)を仮定する．任意の$\pi M\in\qgr A$と任意の下限$N$を取る．
$M$を十分負の添字まで切断しても$\pi M$は変わらないので，
$M_j=0$ ($j>i-4N$)としてよい．
この切断もNoether性により有限生成である．
選択したresidueの成分から生成される部分加群を
\[
 L=MI_i=\sum_q M_{i+4q}A
\]
と置く．$M/L$は有限生成$C_i$加群なので，(2)より有限次元である．
$L$はNoetherianな$M$の部分加群であるから有限生成であり，
有限個の$M_{i-4n}$ ($n\ge N$)の元で生成できる．
よって有限直和からの全射
\[
 \bigoplus_{\alpha}t^{-n_\alpha}E_i\twoheadrightarrow\pi M,
 \qquad n_\alpha\ge N
\]
を得る．これが任意に深い負の軌道による(AZ1)である．
任意の全射のkernelに前補題を適用すれば，十分大きい$n$で
$\Ext^1(t^{-n}E_i,\ker)=0$となり，(AZ2)も得る．

逆に$e_jC_i$が無限次元となる$j$が存在するとする．
$M=e_jC_i$は$P_j$の商なので有限生成で，$\pi M\ne0$である．
すべての$q$で$M_{i+4q}=0$だから，前補題によって
\[
 \Hom(t^{-n}E_i,\pi M)=0\qquad(n\gg0).
\]
十分大きい$N$を選べば，非零対象$t^N\pi M$への
全ての$\Hom(t^{-n}E_i,-)$ ($n\ge0$)が零になる．
これは(AZ1)に反する．最後の主張は
補題~\ref{lem:qa-antiserre-necessary}による．
\end{proof}

\subsection{最小十本矢型の追加矢が与える一般の障害}

\begin{proposition}[最小十本矢型では偶数位相が生成しない]
\label{prop:qa-as10-general-even-obstruction}
任意の$A\in\mathrm{AS}_{10}^{\mathrm{min}}$について，
$(E_i,t)$がAZ ampleであることと$i$が奇数であることは同値である．
特に全ての偶数$i$について
\[
 (\qgr A,E_i)\notin\HH.
\]
この主張に$P$の二捩れ性，avatar曲線の滑らかさは不要である．
\end{proposition}
\begin{proof}
定理~\ref{thm:qa-as10-global-bijection}と
式~\eqref{eq:qa-as10-reduced-potential}により，四周期表示のpotentialは
\[
 W=yq+px+H,
\]
ただし$x:1\to3$, $y:3\to1$は内部次数二，
$p,q$はcycle矢二本の線形結合，$H$はcycle矢四本の線形結合である．
全整数への展開では$x,y$が奇数頂点を二つずつ進める．

全偶数頂点を削除すると，八本のcycle矢はすべて零になり，
奇数頂点を結ぶ$x,y$だけが残る．
全てのcyclic derivativeは，削除されたcycle矢を少なくとも一本含む：
$\partial_xW=p$と$\partial_yW=q$は二本，
cycle矢による微分は一本または三本を含む．
従って残る$x,y$には関係がなく，商は奇数頂点上の一本の有向鎖の
道代数である．その標準右射影$M=e_1C$は，
すべての奇数$j\le1$で$M_j\simeq k$，偶数$j$で$M_j=0$となる
無限次元一元生成加群である．
各偶数位相だけを削除する商はさらに$C$へ全射するので，
前定理の有限次元条件を満たさない．

奇数の一つの位相を削除すると，追加矢$x,y$はともに消え，
残るcycle矢は各周期で長さ高々二の有向鎖になる．
従って削除後の各標準右射影は有限次元であり，前定理を適用できる．
\end{proof}

\begin{remark}[非二捩れでの障害と二周期point module]
この無限path加群はAS10の追加矢から構成され，
中央二重平方表示や中心元$cd+dc$を必要としない．
従って二捩れnoncubicの場合の二周期point moduleによる障害には，
一般$P$でも有効なAS10側の表現がある．
一方，一般avatarに中央二重平方環$B_f$や
通常の次数一生成point schemeが付随すると主張しているのではない．
\end{remark}

\subsection{一般曲線型avatarの全四位相}

以下，$(D,F,G,P)$を本稿のsquare曲線型avatarとする．
$D$は非被約でもよく，$F,G$は指定された$(2,2)$埋め込みの
相異なる二ruling，$P\in\Pic^0(D)\setminus\{\OO_D\}$である．
二つの壁を
\[
 D_F=\{G\simeq FP\},\qquad D_G=\{GP\simeq F\}
\]
と書く．ここで$D_F,D_G$は曲線そのものでなくパラメータ上の条件である．

\begin{theorem}[一般square曲線型avatarの構造頂点]
\label{thm:qa-general-square-pointings}
square AS helixを$(E_i)$とすると，各位相について
強い$\HH$への所属と反Serre軌道のAZ豊富性は同値であり，
次の表で与えられる：
\[
\begin{array}{c|c|c}
 i\bmod4& (E_i,t)\text{ がAZ ample}&(\qgr A,E_i)\in\HH\\ \hline
 0,3&\text{常に成立}&\text{常に成立}\\
 1&GP\not\simeq F&GP\not\simeq F\\
 2&G\not\simeq FP&G\not\simeq FP
\end{array}
\]
\end{theorem}
\begin{proof}
第一mutationの全整数exact比較は
\[
 H_0=\ker(E_0^{\oplus2}\to E_1),\quad
 H_1=E_0,\quad H_2=E_2,\quad H_3=E_3
\]
である．$D_F$上では最小十本矢型なので，前命題によって
$E_2=H_2$の反Serre生成条件が破れる．
$F,G$を交換して第二mutationを行えば，$D_G$上で
$E_1$について同じ結論を得る．

$D_F$の外では定理~\ref{thm:qa-as10-forward-as}により
第一mutationはcubic型AS代数である．この定理は曲線の被約性を要求しない．
命題~\ref{prop:qa-as10-noetherian-qgr}のexact同値は
上記の$H_1=E_0$, $H_2=E_2$, $H_3=E_3$を保つ．
論文A~\cite[\texttt{thm:all-admissible-inverse}]{QuadricIntrinsic2026}は，
\emph{任意の}三次三次元AS正則$\mathbb Z$-代数の標準零頂点を
構造対象とする組が強い$\HH$に属することを述べる．
これを得られたcubic代数および各添字shiftに適用すると，
全$H_i$を構造対象とする組が強い$\HH$に属する．
特に$E_0,E_2,E_3$の所属を得る．
ここで同じrootやcurve上のtensorがnc heartへ延びるとは仮定しない．

第二mutationを用いると，$D_G$の外で$E_0,E_1,E_3$の所属を得る．
二壁の交点では$P^2=\OO_D$, $P\ne\OO_D$, $G=FP$である．
補題~\ref{lem:qa-as10-double-wall-reduced-geometries}により，
この場合$D$はsmooth/nodal/$I_2$の三型に限られる．
全三型についての既存のnoncubic認識と
系~\ref{cor:qa-noncubic-avatar-vertex-h2-status}により
$E_0,E_3$が強い$\HH$に属する．
これで表の全ての肯定・否定を得た．
肯定側のAZ豊富性は補題~\ref{lem:qa-antiserre-necessary}から従い，
否定側はすでに生成障害そのものとして証明されている．
四周期はexact自己同値$t$なので，全整数へ結論が移る．
\end{proof}

\begin{corollary}[Rotated零頂点と一般最小十本矢型のpointing]
\label{cor:qa-general-rotated-as10-pointings}
曲線型rotated avatarを$(D,F_r,H_r,P)$，
$C=H_rF_r^{-2}$とし，標準helixを$(R_i)$とする．
強い$\HH$への所属は
\[
\begin{array}{c|c}
 i\bmod4&(\qgr A_{\rm rsq},R_i)\in\HH\\ \hline
 0&C\not\simeq\OO_D\\
 1&C\not\simeq P^2\\
 2,3&\text{常に成立}
\end{array}
\]
である．とくに零頂点はHilbert曲面が$\Sigma_0$の場合に属し，
$\Sigma_2$の場合には属さない．
また，全ての最小十本矢型avatarでは
構造頂点$H_i$が強い$\HH$に属することと$i$が奇数であることは同値である．
\end{corollary}
\begin{proof}
squareとの辞書は
$F_s=F_r$, $G_s=H_rF_r^{-1}P^{-1}$，
また$R_i=E_{i+1}$である．従って
\[
 G_sP\simeq F_s\iff C\simeq\OO_D,
 \qquad
 G_s\simeq F_sP\iff C\simeq P^2.
\]
前定理に代入すればrotatedの表を得る．
Hilbert曲面の判定は
定理~\ref{thm:qa-core-avatar-square-rotated}による．
この結論は$\coh\Sigma_2$との同一視から得たものではなく，
AS10の無限pathによる生成障害から得ている．

最小十本矢型では$D_F$上にあり，$H_1=E_0$，$H_3=E_3$なので
奇数位相は前定理によって強い$\HH$に属する．
偶数位相は命題~\ref{prop:qa-as10-general-even-obstruction}により除外される．
\end{proof}

\begin{proposition}[古典型の全位相]
\label{prop:qa-classical-all-vertex-pointings}
squareおよびrotatedの各一つの古典型では，
すべての標準頂点を構造対象とする組が強い$\HH$に属する．
\end{proposition}
\begin{proof}
命題~\ref{prop:qa-as-extraction}と
定理~\ref{thm:qa-as-avatar-bijection}の古典型は
$\PP^1\times\PP^1$であり，標準四項はそれぞれ
\[
 (\OO,\OO(1,0),\OO(0,1),\OO(1,1)),\qquad
 (\OO,\OO(-1,1),\OO(0,1),\OO(1,2))
\]
である．全整数helixも線束からなる．各線束とのtensorはexact自己同値で，
$(\coh(\PP^1\times\PP^1),\OO)$は通常のSegre表示による
強い$\HH$の対象だから，結論を得る．
古典型を曲線型の許容条件$P\ne\OO_D$へ形式的に代入したものではない．
\end{proof}

\begin{remark}[範囲と残る比較]
以上は本稿で分類した体上のsquare／rotated／最小十本矢型の
全標準構造頂点を扱う．任意の非標準exceptional対象を構造対象とする問題，
一般avatar族の相対pointed認識，任意基底上のAS/QP群oidの同値は別の課題である．
また，一般のpointed圏について「反Serre ampleなら強い$\HH$」と
主張したのではなく，表の肯定側では実際のcubic/noncubic認識を用いた．
\end{remark}
